\documentclass{beamer}

\usepackage[utf8]{inputenc}

\input{../helper}
\graphicspath{{../img/}}

\title{Introdução à Computação \\ estruturas em C}
\author{Alexandre Rademaker}
\date{}

\begin{document}

\begin{frame}
 \maketitle
\end{frame}

\begin{frame}[fragile,allowframebreaks]{motivação}

  Em um aplicativo de contatos, telefones e nomes não são informações
  separadas, certo?

  Um \textbf{contato} representa uma entrada no banco de dados e cada
  entrada terá \textbf{campos} que descrevem as informações que temos
  sobre o referido contato: nome, telefone, endereço, etc.

  Durante uma busca e/ou ordenação, em geral, queremos podemos
  escolher o campo: por nome, por telefone, por empresa, etc.

  \framebreak

  Vamos imaginar que nossos contatos só tivessem nome e telefone.

  Poderíamos tentar manter os dados em dois diferentes arrays, mas
  como garantir a sincronia? Os elementos em posições iguais nos dois
  arrays correspondem a informações relacionadas ao mesmo contato.

  Manter sincroniza em ambos pode ser difícil.

  \framebreak

  \lstinputlisting[style=tinyc,caption=src/phonebook0.c]{src/phonebook0.c}

\end{frame}  


\begin{frame}[fragile,allowframebreaks]{structuras em C}

  Mas podemos definir nossas próprias estruturas de dados

  \lstinputlisting[style=tinyc,caption=src/phonebook1.c]{src/phonebook1.c}

  \framebreak

  Isto é encapsulamento (encapsulation).

  A idéia de que dados relacionados são agrupados, neste caso, nome e
  número, na mesma estrutura de dados.

  A cor de um pixel, com valores de vermelho, verde e azul, também
  pode ser bem representada por uma estrutura de dados.

  Também podemos imaginar que estruturas para armazenar valores
  racionais ou inteiros grandes. E estruturas podem ter arrays ou
  outras estruturas como \textbf{membros} (campos).

  \href{https://en.wikipedia.org/wiki/Struct_(C_programming_language)}{para saber mais}

  \framebreak

  \begin{lstlisting}[style=normalc]
    struct pointer {
      int x;
      int y;
    } ;
    
    int main(void) {
      struct pointer p; // need the 'struct'
      p.x = 10;
      p.y = 20;
      printf("Point: %d %d\n", p.x, p.y);
    }
  \end{lstlisting}

  A estrutura foi definida e depois uma variável do tipo foi
  declarada.

  \framebreak

  \begin{lstlisting}[style=normalc]
    int main(void) {
      struct {
        int x;
        int y;
      } p;
      
      p.x = 10;
      p.y = 20;
      printf("Point: %d %d\n", p.x, p.y);
    }
  \end{lstlisting}

  A estrutura foi definida e uma variável declarada ao mesmo tipo. Não
  é possível reusar a definição (não foi nomeada).

  \framebreak

  \begin{lstlisting}[style=normalc]
    typedef struct point {
      int x;
      int y;
    } t_point;
    
    int main(void) {
      t_point p;
      p.x = 10;
      p.y = 20;
      printf("Point: %d %d\n", p.x, p.y);
    }
  \end{lstlisting}

  O nome \ilcode{pointer} é opcional, mas necessário para estruturas
  recursivas (mais tarde!)
  
\end{frame}


\end{document}

%%% Local Variables:
%%% mode: latex
%%% TeX-master: t
%%% End:
